\section{Discussion}
The results suggest that the accuracy cost of sparsity in VOS is small once the sparsity pattern is learned jointly with the task rather than fixed a priori, and that the main benefit of cross-modal fusion is not raw accuracy but robustness during occlusion, where a single modality's memory becomes unreliable. Two limitations are worth noting. First, the relevance predictor $g_\theta$ is trained with an auxiliary dense-attention oracle on a subset of steps, which adds training-time overhead that does not appear in our reported inference FPS; a fully sparse training regime is left to future work. Second, our retrieval budget $k$ is fixed across the whole sequence; videos with many simultaneously visible objects may benefit from a budget that scales with the number of active tracks rather than a constant per-query value. We also observed that \methodname{}'s advantage narrows on short clips (under 20 frames), where the memory bank is small enough that dense attention is already cheap; sparsity mainly pays off as sequence length grows, which matches our motivating use case of long, uncut video.
